\documentclass[11pt]{article}
\usepackage[margin=2.2cm]{geometry}
\usepackage[T1]{fontenc}
\usepackage{lmodern}
\usepackage{amsmath,amssymb}
\usepackage{booktabs}
\usepackage{parskip}
\usepackage{fancyvrb}
\usepackage{plumber}
\usepackage[hidelinks]{hyperref}

\DefineVerbatimEnvironment{code}{Verbatim}
  {frame=leftline, framerule=1pt, rulecolor=\color{black!25},
   xleftmargin=1em, fontsize=\small, formatcom=\color{black!80}}
\newenvironment{result}{\par\smallskip\noindent\hspace*{1em}}{\par\medskip}
\newcommand\opt[1]{\texttt{#1}}
\newcommand\cs[1]{\texttt{\textbackslash#1}}
\newcommand\meta[1]{\ensuremath{\langle}\textit{#1}\ensuremath{\rangle}}
\newcommand\marg[1]{\texttt{\{}\meta{#1}\texttt{\}}}
\newcommand\oarg[1]{\texttt{[}\meta{#1}\texttt{]}}

\title{The \textsf{plumber} package\\[2pt]
  \large Diagrams of 2d TQFT cobordisms in \LaTeX{} with Ti\emph{k}Z}
\date{Version 1.0, 2026/10/22}
\author{}

\begin{document}
\maketitle

\noindent
A two-dimensional topological quantum field theory is a symmetric monoidal
functor out of the cobordism category $\mathbf{2Cob}$, whose morphisms are
surfaces between disjoint unions of circles. Every such surface is a
composite of a few generators (cylinders, pairs of pants, caps, \dots), and
this package draws the composite from a description layer by layer:
\[
  \PL{2,3}{pant, id}\PL{3,2}{id, copant}
  \;=\; \PL[reduce]{2,3}{pant, id}\PL{3,2}{id, copant}
\]
Each \cs{PL}\marg{m,n}\marg{generators} is one layer, the tensor product of
the listed generators, going from $m$ to $n$ circles; consecutive \cs{PL}
commands are composed into one surface. The package keeps track of how many
circles there are on every seam, places them so that the tubes run smoothly,
and draws the seams where layers are glued as circles with a solid front and
a dashed back. With \opt{reduce} it computes the topology of the result and
draws the normal form instead.

\section{Quick start}

\begin{code}
\usepackage{plumber}   % loads TikZ
...
\PL{1,2}{pant}\PL{2,2}{swap}\PL{2,1}{copant}
\end{code}
\begin{result}
  \PL{1,2}{pant}\PL{2,2}{swap}\PL{2,1}{copant}
\end{result}

The first argument is the number of input and output circles of the layer
(\verb|2,3| means two circles in, three out); it is checked against the
generators and may also be left empty. The second argument is a
comma-separated list of generators, read from the first circle to the last
(top to bottom in horizontal diagrams, left to right in vertical ones).
\verb|id^3| is short for \verb|id, id, id|, and the same works for any
generator.

\section{Generators}

\begin{center}
\begin{tabular}{@{}lllcc@{}}
  \toprule
  name & aliases & circles & $\chi$ & symbol \\
  \midrule
  \opt{id}     & \opt{cyl}, \opt{cylinder} & $1\to1$ & $0$  & $\mathrm{id}$ \\
  \opt{pant}   & \opt{pants}              & $1\to2$ & $-1$ & $\Delta$ \\
  \opt{copant} & \opt{copants}            & $2\to1$ & $-1$ & $\mu$ \\
  \opt{unit}   & \opt{birth}              & $0\to1$ & $1$  & $\eta$ \\
  \opt{counit} & \opt{death}              & $1\to0$ & $1$  & $\varepsilon$ \\
  \opt{swap}   & \opt{twist}              & $2\to2$ & $0$  & $\sigma$ \\
  \opt{iswap}  & \opt{itwist}             & $2\to2$ & $0$  & $\sigma^{-1}$ \\
  \opt{pair}   & \opt{ev}                 & $2\to0$ & $0$  & $\beta$ \\
  \opt{copair} & \opt{coev}               & $0\to2$ & $0$  & $\gamma$ \\
  \bottomrule
\end{tabular}
\end{center}

Diagrams are read in the direction of composition: from left to right, or
from top to bottom with \opt{vertical}. So \opt{pant} splits one circle into
two and \opt{copant} merges two circles into one. \opt{swap} and \opt{iswap}
cross the two tubes with the second (respectively the first) tube in front;
as cobordisms they are equal. \opt{pair} and \opt{copair} are the bent
cylinders (U-turns) that make $\mathbf{2Cob}$ a rigid category.
\cs{PLgenerators} draws them all:

\begin{code}
\PLgenerators[names]
\PLgenerators*[vertical, names, list={pant, copant, unit, counit, iswap}]
\end{code}
\begin{result}
  \PLgenerators[names]
\end{result}
\begin{result}
  \PLgenerators*[vertical, names, list={pant, copant, unit, counit, iswap}]
\end{result}

\cs{PLgenerators}\oarg{options} draws the generators in the key \opt{list}
(by default all of them except \opt{iswap}) next to each other; the starred
form leaves out \opt{id}.

\section{The features}

\subsection{Composition}

Consecutive \cs{PL} commands form one diagram (spaces and single line breaks
between them are allowed; an empty line or any other token ends the
diagram). The output circles of one layer are glued to the input circles of
the next one, so the counts must agree. Options can be given to any \cs{PL}
of the chain and apply to the whole diagram.

\begin{code}
\[ \PL{1,2}{pant}\PL{2,3}{pant, id} = \PL{1,2}{pant}\PL{2,3}{id, pant}
   \qquad
   \PL{1,2}{id, unit}\PL{2,1}{copant} = \PL{1,1}{id} \]
\end{code}
\begin{result}
  \( \PL{1,2}{pant}\PL{2,3}{pant, id} = \PL{1,2}{pant}\PL{2,3}{id, pant}
     \qquad
     \PL{1,2}{id, unit}\PL{2,1}{copant} = \PL{1,1}{id} \)
\end{result}

\begin{code}
\[ \PL{1,3}{id, copair}\PL{3,1}{pair, id} = \PL{1,1}{id}
   \qquad
   \PL{2,2}{id^2}\PL{2,4}{pant, pant}\PL{4,4}{id, swap, id}\PL{4,2}{copant^2} \]
\end{code}
\begin{result}
  \( \PL{1,3}{id, copair}\PL{3,1}{pair, id} = \PL{1,1}{id}
     \qquad
     \PL{2,2}{id^2}\PL{2,4}{pant, pant}\PL{4,4}{id, swap, id}\PL{4,2}{copant^2} \)
\end{result}

If the circles do not match, or a generator is unknown, the package reports
an error that names the offending layer and draws a red {\color{red}\bfseries
??} instead of the diagram. Should \cs{PL} already be defined by another
package, \cs{PLlayer} is a synonym with the same arguments (the two can be
mixed in a chain).

\subsection{Horizontal or vertical}

Diagrams are \opt{horizontal} by default: composition runs from left to
right and the first circle is on top. With \opt{vertical}, composition runs
from top to bottom and the first circle is on the left. \opt{reverse} flips
the direction of composition (right to left, or bottom to top). The circles
at the start of the diagram are seen from the outside and drawn whole; all
other seams show their front half solid and their back half dashed.
\opt{orientation=horizontal} and \opt{orientation=vertical} are long forms,
and \cs{plset}\marg{options} changes the defaults for the rest of the group.

\begin{code}
\PL[vertical]{2,3}{pant, id}\PL{3,2}{id, copant}
\PL[vertical, reverse]{2,3}{pant, id}\PL{3,2}{id, copant}
\PL[reverse]{2,3}{pant, id}\PL{3,2}{id, copant}
\end{code}
\begin{result}
  \PL[vertical]{2,3}{pant, id}\PL{3,2}{id, copant}\qquad
  \PL[vertical, reverse]{2,3}{pant, id}\PL{3,2}{id, copant}\qquad
  \PL[reverse]{2,3}{pant, id}\PL{3,2}{id, copant}
\end{result}

\subsection{Reduced form}

Up to diffeomorphism, a cobordism is determined by which boundary circles
lie on the same connected component and by the genus of every component.
With \opt{reduce}, the package computes these invariants (by tracking the
glued circles and adding up Euler characteristics) and draws the normal form:
first the input circles are permuted so that every component gets
neighbouring circles, then each component merges its inputs with a tree of
\opt{copant}s, gets its handles as \opt{pant}/\opt{copant} pairs, and splits
into its outputs with a tree of \opt{pant}s, and finally the outputs are
permuted into place. Components with no inputs start with a \opt{unit},
components without outputs end with a \opt{counit}, and the two kinds of
cylinders that bend back (two inputs or two outputs, genus~0) become
\opt{pair} and \opt{copair}.

\begin{code}
\[ \PL{1,3}{copair, id}\PL{3,3}{swap, id}\PL{3,2}{id, copant}
   = \PL[reduce]{1,3}{copair, id}\PL{3,3}{swap, id}\PL{3,2}{id, copant} \]
\[ \PL{2,4}{pant, pant}\PL{4,4}{id, swap, id}\PL{4,2}{copant, copant}
   = \PL[reduce]{2,4}{pant, pant}\PL{4,4}{id, swap, id}\PL{4,2}{copant, copant} \]
\end{code}
\begin{result}
  \( \PL{1,3}{copair, id}\PL{3,3}{swap, id}\PL{3,2}{id, copant}
     = \PL[reduce]{1,3}{copair, id}\PL{3,3}{swap, id}\PL{3,2}{id, copant} \)
\end{result}
\begin{result}
  \( \PL{2,4}{pant, pant}\PL{4,4}{id, swap, id}\PL{4,2}{copant, copant}
     = \PL[reduce]{2,4}{pant, pant}\PL{4,4}{id, swap, id}\PL{4,2}{copant, copant} \)
\end{result}

Closed components (surfaces without boundary) are drawn as separate factors,
a sphere or a surface with handles, before the rest. With
\opt{closed=factor} they are written as numbers instead: a closed surface of
genus $g$ evaluates to $Z(\Sigma_g)$ in a TQFT $Z$, and the key
\opt{closed value} sets this expression (\verb|#1| is the genus). The
\opt{separator} key puts a symbol between the factors.

\begin{code}
\[ \PL{0,1}{unit}\PL{1,2}{pant}\PL{2,2}{swap}\PL{2,1}{copant}\PL{1,0}{counit}
   = \PL[reduce]{0,1}{unit}\PL{1,2}{pant}\PL{2,2}{swap}\PL{2,1}{copant}\PL{1,0}{counit} \]
\[ \PL{0,3}{unit, copair}\PL{3,1}{id, pair}\PL{1,0}{counit}
   = \PL[reduce, separator=\sqcup]{0,3}{unit, copair}\PL{3,1}{id, pair}\PL{1,0}{counit}
   = \PL[reduce, closed=factor]{0,3}{unit, copair}\PL{3,1}{id, pair}\PL{1,0}{counit} \]
\[ \PL{1,2}{id, unit}\PL{2,3}{id, pant}\PL{3,2}{id, copant}\PL{2,1}{id, counit}
   = \PL[reduce, closed=factor, closed value={\lambda_{#1}}]
       {1,2}{id, unit}\PL{2,3}{id, pant}\PL{3,2}{id, copant}\PL{2,1}{id, counit} \]
\end{code}
\begin{result}
  \( \PL{0,1}{unit}\PL{1,2}{pant}\PL{2,2}{swap}\PL{2,1}{copant}\PL{1,0}{counit}
     = \PL[reduce]{0,1}{unit}\PL{1,2}{pant}\PL{2,2}{swap}\PL{2,1}{copant}\PL{1,0}{counit} \)
\end{result}
\begin{result}
  \( \PL{0,3}{unit, copair}\PL{3,1}{id, pair}\PL{1,0}{counit}
     = \PL[reduce, separator=\sqcup]{0,3}{unit, copair}\PL{3,1}{id, pair}\PL{1,0}{counit}
     = \PL[reduce, closed=factor]{0,3}{unit, copair}\PL{3,1}{id, pair}\PL{1,0}{counit} \)
\end{result}
\begin{result}
  \( \PL{1,2}{id, unit}\PL{2,3}{id, pant}\PL{3,2}{id, copant}\PL{2,1}{id, counit}
     = \PL[reduce, closed=factor, closed value={\lambda_{#1}}]
         {1,2}{id, unit}\PL{2,3}{id, pant}\PL{3,2}{id, copant}\PL{2,1}{id, counit} \)
\end{result}

After every diagram (reduced or not) the following macros expand to the
invariants of the surface that was just drawn: \cs{PLlastcomponents} (number
of connected components), \cs{PLlastclosed} (how many of them are closed),
\cs{PLlastgenus} (sum of the genera), \cs{PLlasteuler} (Euler
characteristic) and \cs{PLlastnormalform} (the normal form as a string of
\cs{PL} commands, closed factors separated by \verb|\cdot|). For the last
example: \PLlastcomponents{} components, \PLlastclosed{} closed, genus
\PLlastgenus, $\chi = \PLlasteuler$, and the normal form is
\begin{center}\ttfamily\small\PLlastnormalform\end{center}

\subsection{Names}

The key \opt{names} writes the symbol of each layer under it (or to its
right in vertical diagrams). The symbols are set with \opt{symbols}, which
takes a key--value list of generator names and math material.

\begin{code}
\PL[names]{2,3}{pant, id}\PL{3,2}{id, copant}
\PL[vertical, names, symbols={pant=m^*, copant=m}]{2,3}{pant, id}\PL{3,2}{id, copant}
\end{code}
\begin{result}
  \PL[names]{2,3}{pant, id}\PL{3,2}{id, copant}\qquad
  \PL[vertical, names, symbols={pant=m^*, copant=m}]{2,3}{pant, id}\PL{3,2}{id, copant}
\end{result}

\section{Appearance}

\begin{center}
\begin{tabular}{@{}lll@{}}
  \toprule
  key & default & meaning \\
  \midrule
  \opt{circle sep}  & \opt{1}     & distance between neighbouring circles (cm) \\
  \opt{radius}      & \opt{0.3}   & radius of a circle (cm) \\
  \opt{length}      & \opt{1.2}   & length of a layer (cm) \\
  \opt{perspective} & \opt{0.3}   & width of the ellipses relative to their height \\
  \opt{crotch}      & \opt{0.6}   & how far the crotch of a pair of pants reaches back \\
  \opt{gap}         & \opt{0.4}   & space between the factors of a reduced diagram (cm) \\
  \opt{seams}       & \opt{true}  & draw the circles where layers are glued \\
  \opt{smooth}      & \opt{4}     & maximal number of layout passes \\
  \opt{scale}       & \opt{1}     & scales the picture (not the text) \\
  \opt{baseline}    & \opt{-0.5ex}& vertical alignment, as in TikZ \\
  \opt{tikz}        &             & further options for the \texttt{tikzpicture} \\
  \opt{fill}        & \opt{white} & colour of the surface \\
  \bottomrule
\end{tabular}
\end{center}

The drawing uses the TikZ styles \opt{pl surface} (the white area of every
piece), \opt{pl outline}, \opt{pl boundary} (whole circles), \opt{pl seam}
(front half of a seam), \opt{pl hidden} (its dashed back half), \opt{pl name},
\opt{pl name right}, \opt{pl separator}, \opt{pl coefficient} and
\opt{pl error}. The keys \opt{surface style}, \opt{outline style},
\opt{boundary style}, \opt{seam style} and \opt{hidden style} append to the
corresponding styles.

\begin{code}
\PL[fill=blue!10, length=0.8, circle sep=0.8, radius=0.25]{1,2}{pant}\PL{2,1}{copant}
\PL[seams=false, outline style={thick, blue!60!black}]{1,2}{pant}\PL{2,1}{copant}
\PL[hidden style={densely dotted}, scale=0.7]{1,2}{pant}\PL{2,1}{copant}
\end{code}
\begin{result}
  \PL[fill=blue!10, length=0.8, circle sep=0.8, radius=0.25]{1,2}{pant}\PL{2,1}{copant}\qquad
  \PL[seams=false, outline style={thick, blue!60!black}]{1,2}{pant}\PL{2,1}{copant}\qquad
  \PL[hidden style={densely dotted}, scale=0.7]{1,2}{pant}\PL{2,1}{copant}
\end{result}

\section{Layout}

The numbers of circles change from seam to seam, so their positions are
computed. Every generator says where it would like the circles on the
neighbouring seams to be: a cylinder keeps the position, a pair of pants puts
its legs half a circle distance on either side of the waist, a swap exchanges
its two positions, and generators without inputs (\opt{unit}, \opt{copair})
attach to their neighbours with a weak preference. On every seam the
positions are the best least-squares fit to these wishes under the condition
that neighbouring circles are at least \opt{circle sep} apart (an isotonic
regression), and a few sweeps back and forth through the diagram let both
neighbours of a seam have their say. The result is the familiar picture:
straight cylinders stay straight and pairs of pants are symmetric whenever
possible.

\section{The standalone class}

With \verb|\documentclass[tikz]{standalone}| every TikZ picture becomes a
page of its own. Diagrams inside \verb|\[...\]| and \verb|\(...\)| are
handled like in the \textsf{temperleylieb} package: the whole formula becomes
one cropped page. Diagrams outside formulas are pages of their own. Other
math environments (\verb|$...$|, \verb|equation|) cannot be turned into
pages; the package warns and draws the diagram inline, so use
\verb|\[...\]| there, or the class option \opt{varwidth}.

\section{Summary of commands}

\begin{description}
  \item[\cs{PL}\oarg{options}\marg{m,n}\marg{generators}] One layer; a
    chain of layers is one diagram. \cs{PLlayer} is a synonym.
  \item[\cs{PLgenerators}\texttt{*}\oarg{options}] The generators in
    \opt{list} side by side.
  \item[\cs{plset}\marg{options}] Sets options for the current group.
  \item[\cs{PLlastcomponents}, \cs{PLlastclosed}, \cs{PLlastgenus},
    \cs{PLlasteuler}] Invariants of the last diagram.
  \item[\cs{PLlastnormalform}] Normal form of the last diagram.
\end{description}

The package is written in \textsf{expl3} and needs a \LaTeX{} kernel from
2020-10 or newer and TikZ. It is distributed under the \LaTeX{} Project
Public License~1.3c.

\end{document}
